\documentclass[10pt,a4paper]{amsart}
\usepackage[margin=19mm]{geometry}
\usepackage{amsmath,amssymb,mathtools}
\usepackage[scr=rsfs,cal=euler]{mathalfa}
\usepackage{xfrac}
\usepackage[unicode,hidelinks,bookmarks=false]{hyperref}
\newcommand{\ie}{i.e.\ }
\newcommand{\cf}{cf.\ }
\newcommand{\resp}{respectively\ }
\newcommand{\Subsection}{\subsection}
\providecommand{\nobreakdash}{\nobreak\hbox{-}}
\setlength{\emergencystretch}{2em}
\multlinegap0pt
\usepackage{tikz}
\usepackage{amssymb}
\usepackage[shortlabels]{enumitem}
\usepackage{float}
\newtheorem{theorem}{Theorem}
\newtheorem{lemma}{Lemma}
\newcommand{\epn}{\mbox{epn}}
\newcommand{\gammaenaL}{\gamma_1^{\mathrm L}}
\newcommand{\gammakL}{\gamma_k^{\mathrm L}}
\title{On $k$-limited domination in graphs}
\author{Dragana Bo\v{z}ovi\'c, Gordana Radi\'c, \v{Z}ana Kovijani\'c-Vuki\'cevi\'c, Aleksandra Tepeh}
\date{Source: arXiv:2606.22422v1 (CC BY 4.0)}
\begin{document}
\maketitle

\noindent\emph{Source and licence.} On $k$-limited domination in graphs. Version arXiv:2606.22422v1 (CC BY 4.0), distributed by arXiv under the Creative Commons Attribution 4.0 International licence, http://creativecommons.org/licenses/by/4.0/, as recorded in the arXiv licence metadata for this exact version and shown on the abstract page https://arxiv.org/abs/2606.22422v1. Full licence notice: \texttt{licenses/arxiv-2606.22422-CC-BY-4.0.txt}.\par\smallskip

\noindent\textit{Editorial scope.} Counted unit: the article's theorem zg-kar (source0.tex lines 823--831). The article's complete original proof of it is delivered with it (source0.tex lines 834--865, one \texttt{proof} environment of 32 lines). No mathematics is written, repaired or re-derived: the statement and the proof are the article's own lines, in the article's own notation, and no retained line is substituted or re-flowed.  Retained uncounted as the source-local context the proof needs: lines 304--312; lines 688--690; lines 754--758.  Nothing was omitted from any retained span, so the package carries no omission record.  No retained line contains a citation, so the wrapper carries no bibliography and no bibliography entry.  No retained line contains a figure environment, an image-inclusion command or drawing code, so no image is delivered and the package declares no asset dependency.  Editorial formatting only, in addition: an AMS \texttt{amsart} preamble that restates the article's own declarations and macros used by the retained lines, namely the environments \texttt{theorem}, \texttt{lemma} on separate counters, together with the standard packages those lines need.  The retained environments print in this document as Theorem 1, Lemma 1 in order of appearance, whereas the article prints the same items with the numbering of its own full document.  The complete native source of the article as distributed in the arXiv e-print for this exact version is bundled unchanged as \texttt{sources/203345/source0.tex}.
\begin{theorem} 
\label{bounds}
Let $G$ be a connected graph of order $n\geq 3$, and $k < \Delta(G)$. Then
$$
\max\left\{\gamma(G),\ \left\lceil \frac{n}{k+1} \right\rceil\right\} 
\leq \gammakL(G)  \leq 
n-k.
$$
\end{theorem}

\begin{lemma}\label{k=1:universal}
Let $G$ be a connected graph on $n\geq 3$ vertices containing a universal vertex. Then $\gammaenaL(G)=n-1$.  
\end{lemma}

\begin{lemma}\label{pendant}
Let $G$ be a connected graph on $n\geq 3$ vertices. 
If $\gamma^{\mathrm L}_1(G)=n-1$, then $\mathrm{diam}(G)\le 2$,
and every non-pendant edge of $G$ lies in a triangle or $G$ contains a universal vertex.
\end{lemma}

\begin{theorem}\label{zg-kar}
Let $G$ be a connected graph on $n\geq 3$ vertices. Then $\gammaenaL(G)=n-1$ if and only if
\begin{enumerate}
\item[(i)] $\mathrm{diam}(G)\leq 2$,
\item[(ii)] every non-pendant edge of $G$ lies in a triangle, or $G$ has a universal vertex, and
\item[(iii)] $G$ contains no subset $S \subseteq V(G)$ with $|S|\ge 3$
such that $\epn(s,S)\neq\emptyset$ for every $s\in S$, and $N(S)=\displaystyle \bigcup _{s\in S}\epn(s,S)$.
\end{enumerate}
\end{theorem}

\begin{proof}
Suppose first that $\gammaenaL(G)=n-1$. 
Then conditions $(i)$ and $(ii)$ follow from Lemma \ref{pendant}. To prove condition $(iii)$, suppose to the contrary, that $G$ contains a subset $S \subseteq V(G)$ with $|S|\ge 3$
such that $\epn(s,S)\neq\emptyset$ for every $s\in S$, and $N(S)= \bigcup _{s\in S}\epn(s,S)$.
Clearly, $\epn(s,S)\cap \epn(t,S)=\emptyset$ for any different vertices $s$ and $t$ from $S$. Thus each vertex in $N(S)$ is adjacent to precisely one vertex in $S$. 
Now, let $D=V(G)-S$.
Since $\epn(s,S)\neq\emptyset$ for every $s\in S$, $s$ has a neighbor in $D$, and so $D$ is a dominating set. Moreover, each vertex $d\in D$, that has a neighbor in $S$, lies in $N(S)$ and therefore
has exactly one neighbor in $S=V(G)-D$. 
Thus $D$ is a $1$-limited dominating set, and $\gammaenaL(G)\leq |D|=n-|S|\leq n-3$, a contradiction.




To prove the converse, assume that $G$ satisfies conditions $(i)-(iii)$, and let $D$ be a $\gammaenaL(G)$-set. 
If $G$ has a universal vertex, the claim follows from Lemma \ref{k=1:universal}. Hence in the sequel we assume, that $G$ has no universal vertex. 

Suppose that $|D|\leq n-2$, and let $S=V(G)-D$.
Then there exist distinct vertices $s_1, s_2\in S$. Since $\text{diam}(G)\leq 2$, either $s_1s_2\in E(G)$ or $d(s_1,s_2)=2$.  
First, let $s_1s_2\in E(G)$. If $\deg(s_i)=1$ for some $i\in\{1,2\}$, then $s_i$ does not have a neighbor in $D$, a contradiction. Hence, $s_1s_2$ is a non-pendant edge, and by the condition $(ii)$ there exists a vertex $s_3\in V(G)$ adjacent to both $s_1$ and $s_2$. Clearly $s_3\notin D$, since otherwise $s_3$ would dominate two vertices outside $D$, violating the $1$-limited constraint. 
In the second case, when $d(s_1,s_2)=2$, there exists $s_3\in V(G)$ such that $s_1s_3s_2$ is a path of length $2$, immediately implying that $s_3\in S$. Thus $|S|\geq 3$. 
%Now we show that $S$ violates condition $(iii)$.

Since $D$ is a dominating set,
every vertex $s\in S$ has at least one neighbor in $D$, say  $d$, and since $D$ is $1$-limited, $d$ has at most one neighbor
in $S$. Thus $N(d)\cap S=\{s\}$ and so $d\in \epn(s,S)$, implying $\epn(s,S)\neq\emptyset$
for every $s\in S$.
Moreover, it is straightforward to see that the sets 
$N(S)=\{d\in D ; N(d)\cap S\neq\emptyset\}$ and $ \bigcup_{s\in S}\epn(s,S)$ are equal.
But since $|S|\ge 3$, we now have a contradiction with condition $(iii)$.

Consequently, $|D| \ge n - 1$, while the reverse inequality is guaranteed by Theorem~\ref{bounds}, yielding $\gamma^{\mathrm L}_1(G) = n - 1$.
\end{proof}

\end{document}
